%\vspace{-2mm}
\section{Related Work} \label{sec:relwork}
%\vspace{-3mm}
LSR and PCR are fundamental tasks in legal document processing and several research works have been done in this area, for example, techniques based on n-grams \cite{salton1988termweight,zeng2007knowledge}, doc embeddings \cite{le2014doc2vec}, transformers \cite{vold2021tflegal} and others \cite{salton1975vector,robertson2009bm25,liu2023ml,hofmann2013balancing,wang2018modeling,ma2022incorporating}. 
Various works have been done for identifying legal statutes~\cite{wang2018modeling,wang2019hmn,chalkidis2019neural,zhong2018topjudge,wu2023precedent,paul2024legal}. 
Similarly, various works have focused on prior case retrieval~\cite{rabelo2022semantic,joshi2023ucreat,tang2024casegnn,tang2024caselink,qin2024ljpretrieval,bhattacharya2020hierspcnet,ma2021retrieving,ma2024llmrelevance}.
% Due to space limitations, we provide more details about above-mentioned works in App. \ref{app:related-work}.
More details regarding related work are provided in App. \ref{app:related-work}.
% To the best of our knowledge, the two tasks have only been attempted independently. In this paper, we are the first ones to provide a parallel corpus that provides an opportunity to develop models that exploit the dependence between the two tasks. 
To our knowledge, the two tasks have only been addressed separately. We present the first parallel corpus enabling models to leverage their interdependence.
%a common model architecture for both the tasks, since both the tasks are legal document retrieval tasks. 

% Identifying the legal statutes and relevant prior cases given a legal fact or situation is one of the most fundamental tasks in law.
% Traditionally, researchers have used statistical and lexical approaches to solve both tasks independently. The advent of deep learning NLP approaches has led to renewed efforts in both tasks using advanced architectures.

% \subsection{Overview of Prior Works} \label{sec:relwork/overview}

% Traditional approaches for identifying relevant statutes and precedents have mostly involved exploiting lexical features such as n-grams of words~\citep{salton1988termweight}, hand-crafted features~\citep{zeng2007knowledge} or embeddings from pre-trained models like Doc2vec~\citep{le2014doc2vec}.
% Lately, transformer-based embedding methods have been used for directly calculating dot product scores between the query and statute/precedent~\citep{vold2021tflegal}.
% While most unsupervised approaches have utilized methods like Vector Space Model~\citep{salton1975vector} and BM25~\citep{robertson2009bm25}, supervised approaches for both tasks can broadly be divided into classification~\citep{liu2023ml,hofmann2013balancing} (model predicts similarity between query and statute/precedent) and ranking based~\citep{wang2018modeling,ma2022incorporating} (model ranks a list of statutes/precedents based on relevance to the query) approaches. 

% \subsection{Identifying Legal Statutes} \label{sec:relwork/statutes}
% Historically, researchers have used multi-label learning frameworks to identify relevant statutes for a query~\citep{wang2018modeling,wang2019hmn,chalkidis2019neural}. In many jurisdictions, identifying the relevant statutes is considered to be a subtask of the broader task of Legal Judgment Prediction~\citep{zhong2018topjudge}, which could entail predicting the legal charges and term of punishment as auxiliary tasks.
% Some approaches have only considered the text of the queries in the classification pipeline, relying on the encoder to generate good quality representations of the query~\citep{chalkidis2019neural}.
% Others have incorporated the text of the statutes as well, in generating statute-aware query representations which are then used for classification~\citep{wang2018modeling,wang2019hmn}.
% It should be noted that most of these approaches have worked in a setup with limited number of statutes (<200), and hence the classification approach suffices.
% Lately, LLMs have been used to perform the task of statute identification~\citep{wu2023precedent}, and these models can utilize their superior language understanding capabilities as well as knowledge of legal statutes to excel in the task of statute identification.

% \subsection{Identifying Precedents} \label{sec:relwork/precedents}
% Unlike statutes, most prior works on prior case retrieval have modeled the task in a ranking framework. The major challenge in this task is the fact that both the queries and precedents are very long. 
% Additionally, it has been observed that the query consists of several legal aspects, and each individual aspect leads to matching with certain precedents that eventually get cited~\citep{rabelo2022semantic}. 
% Mostly, researchers have tried to reduce the noise in the query text by using event information~\citep{joshi2023ucreat,tang2024casegnn}, or extracting salient portions of the document~\citep{qin2024ljpretrieval}. 
% \citet{rabelo2022semantic} took a granular approach, by dividing both the queries and precedents into paragraphs/sentences, scoring each pair of query and precedent sentence, and then generating aggregate scores.
% Other approaches have involved usage of GNNs~\citep{tang2024casegnn,tang2024caselink}, citation network structures~\citep{bhattacharya2020hierspcnet}, making use of the statutes cited from the precedent cases (not the queries)~\citep{qin2024ljpretrieval}, and re-ranking approaches based on some first stage retriever like BM25~\citep{ma2021retrieving}.
% LLMs have also been lately used to summarize the queries and precedents~\citep{qin2024ljpretrieval}, or perform query expansion based on its inherent domain knowledge~\citep{ma2024llmrelevance}.